\documentclass[12pt,a4paper]{article}
\usepackage{cmap} % Makes the PDF copiable. See http://tex.stackexchange.com/a/64198/25761
\usepackage[T1]{fontenc}
\usepackage[brazil]{babel}
\usepackage[utf8]{inputenc}
\usepackage{amsmath}
\usepackage{amsfonts}
\usepackage{amssymb}
\usepackage{amsthm}
\usepackage{textcomp} % \degree
\usepackage{gensymb} % \degree
\usepackage[usenames,svgnames,dvipsnames]{xcolor}
\usepackage{hyperref}
\usepackage{graphicx}
\usepackage[margin=2cm]{geometry}

\hypersetup{
    colorlinks = true,
    allcolors = {blue}
}

\newcommand*\diff{\mathop{}\!\mathrm{d}}
\newcommand*\sen{\operatorname{sen}}

\newcommand*\tipo{PROVA II}
\newcommand*\disciplina{CDI0001}
\newcommand*\eu{Helder G. G. de Lima}
\newcommand*\data{16/10/2015}

\author{\eu}
\title{\tipo - \disciplina}
\date{\data}

\begin{document}
\thispagestyle{empty}
\newgeometry{margin=2cm,bottom=0.5cm}
\begin{center}
\includegraphics[width=9.0cm]{marca}
\noindent\begin{tabular}{l c c r}
  \textbf{\disciplina (turma: \rule{4cm}{0.01cm})}
& \textbf{\tipo}
& \textbf{\data}
\end{tabular}
\\ Prof. \eu\footnote{
Este é um material de acesso livre distribuído sob os termos da licença \href{https://creativecommons.org/licenses/by-sa/4.0/deed.pt_BR}{Creative Commons BY-SA 4.0}.}
\end{center}

\noindent Nome do(a) aluno(a): \rule{13cm}{0.01cm}

%\section*{Instruções}
\begin{center}\fbox{
\begin{minipage}{14cm}

{\footnotesize
\begin{itemize}
\renewcommand{\theenumi}{\Roman{enumi}}
\item Identifique-se em todas as folhas.
\item Mantenha o celular e os demais equipamentos eletrônicos desligados durante a prova.
\item Justifique cada resposta com cálculos ou argumentos baseados na teoria estudada.
\item Resolva (integralmente) apenas os itens de que precisar para somar 10,0 pontos.
\end{itemize}
}

\end{minipage}
}
\end{center}


%\section*{Questões}

\begin{enumerate}
\item (0,5) Dê exemplos de funções deriváveis $f, g: \mathbb{R} \to \mathbb{R}$ tais que $\left[\frac{f(x)}{g(x)} \right]^\prime \neq \frac{f^\prime(x) }{ g^\prime(x) }$. Faça as contas para mostrar que realmente são diferentes.

\item (1,0) Verifique que se $f$, $g$ e $h$ são deriváveis então
\[
\frac{ \diff{} }{\diff{x}}\left[f(x) g(x) h(x) \right]
=
  \frac{ \diff{f(x)} }{\diff{x}} g(x) h(x)
+ f(x) \frac{\diff{g(x)} }{\diff{x}} h(x)
+ f(x) g(x) \frac{\diff{h(x)} }{\diff{x}}.
\]
\item (0,5) Calcule $h^\prime(0)$, sendo $h: \mathbb{R} \to \mathbb{R}$ definida por $h(x) = \sen(x) \cos(x) e^x$.

\item (1,5) Mostre que se $y = f(x) = \pi e^{-2x} + 4e^{3x}$ então vale $y^{\prime \prime} - y^\prime = 6 y$, sendo $y^\prime$ e $y^{\prime\prime}$, respectivamente, as derivadas de primeira e segunda ordem da função indicada.

\item (1,5) Utilize uma aproximação linear para estimar o valor de $\sqrt[3]{8,02}$.

\item (2,0) Ao esvaziar uma bola de futebol, o seu volume (em $cm^3$) após $t$ segundos é dado por $V(t) = 400 \pi (10 - t)$. Com que velocidade ($cm/s$) o raio diminui quando ele mede $5cm$?

\item (2,5) Determine a equação da reta tangente ao gráfico de $f(x) = \dfrac{4}{x + 1}$ em $x = 1$.

\item (2,5) Em que pontos $(x,y)$ da curva $ x^\frac{2}{3} + y^\frac{2}{3} = 8$ a reta tangente é paralela à reta $y = x + 4$?

\end{enumerate}

\newpage
\restoregeometry
\section*{Respostas e observações}


\begin{enumerate}
\item \textit{Dê exemplos de funções deriváveis $f, g: \mathbb{R} \to \mathbb{R}$ tais que $\left[\frac{f(x)}{g(x)} \right]^\prime \neq \frac{f^\prime(x) }{ g^\prime(x) }$. Faça as contas para mostrar que realmente são diferentes.}
Há inúmeros casos em que isso acontece. Por exemplo,
\begin{enumerate}
\item Se $f(x) = 1$ e $g(x) = x$ então
\[
\left[\frac{f(x)}{g(x)} \right]^\prime
= \left[\frac{1}{x} \right]^\prime
= \frac{-1}{x^2}
\neq 0
= \frac{0 }{ 1 }
= \frac{f^\prime(x) }{ g^\prime(x) }
\]

\item Se $f(x) = e^x = g(x)$ então
\[
\left[\frac{f(x)}{g(x)} \right]^\prime
= \left[\frac{e^{2x}}{e^x} \right]^\prime
= \left[ e^x \right]^\prime
= e^x
\neq 1
= \frac{e^x }{ e^x }
= \frac{f^\prime(x) }{ g^\prime(x) }
\]


\end{enumerate}

\item \textit{Verifique que se $f$, $g$ e $h$ são deriváveis então}
\[
\frac{ \diff{} }{\diff{x}}\left[f(x) g(x) h(x) \right]
=
  \frac{ \diff{f(x)} }{\diff{x}} g(x) h(x)
+ f(x) \frac{\diff{g(x)} }{\diff{x}} h(x)
+ f(x) g(x) \frac{\diff{h(x)} }{\diff{x}}.
\]

Basta observar que $f(x) g(x) h(x)$ é o produto da função $f(x) g(x)$ com a função $h(x)$, e aplicar a regra do produto como de costume:
\begin{align*}
\frac{ \diff{} }{\diff{x}}\left[f(x) g(x) h(x) \right]
& = \frac{ \diff{} }{\diff{x}}\left[(f(x) g(x)) \cdot h(x) \right] \\
& =
\left( \frac{ \diff{f(x)} }{\diff{x}} g(x)
+ f(x) \frac{\diff{g(x)} }{\diff{x}} \right) h(x)
+ f(x) g(x) \frac{\diff{h(x)} }{\diff{x}} \\
& =
  \frac{ \diff{f(x)} }{\diff{x}} g(x) h(x)
+ f(x) \frac{\diff{g(x)} }{\diff{x}} h(x)
+ f(x) g(x) \frac{\diff{h(x)} }{\diff{x}}.
\end{align*}

\item \textit{Calcule $h^\prime(0)$, sendo $h: \mathbb{R} \to \mathbb{R}$ definida por $h(x) = \sen(x) \cos(x) e^x$.}

Pelo item anterior, tem-se
\begin{align*}
h^\prime(x)
& = (\sen(x))^\prime \cos(x) e^x + \sen(x) (\cos(x))^\prime e^x + \sen(x) \cos(x) (e^x)^\prime \\
& = \cos(x) \cos(x) e^x + \sen(x) (-\sen(x)) e^x + \sen(x) \cos(x) e^x \\
& = (\cos^2(x) - \sen^2(x) + \sen(x) \cos(x)) e^x
\end{align*}
Em particular, $h^\prime(0) = (\cos^2(0) - \sen^2(0) + \sen(0) \cos(0)) e^0 = (1 - 0 + 0 \cdot 1) 1 = 1.$

\item \textit{Mostre que se $y = f(x) = \pi e^{-2x} + 4e^{3x}$ então vale $y^{\prime \prime} - y^\prime = 6 y$, sendo $y^\prime$ e $y^{\prime\prime}$, respectivamente, as derivadas de primeira e segunda ordem da função indicada.}

Observe que $y^\prime = (\pi e^{-2x} + 4e^{3x})^\prime = \pi e^{-2x} \cdot (-2) + 4 e^{3x} \cdot 3 = -2 \pi e^{-2x} + 12 e^{3x}$ e que então $y^{\prime\prime} = (-2 \pi e^{-2x} + 12 e^{3x})^\prime = 4 \pi e^{-2x} + 36 e^{3x}$. Logo,
\[
y^{\prime \prime} - y^\prime
= (4 \pi e^{-2x} + 36 e^{3x})-(-2 \pi e^{-2x} + 12 e^{3x})
= 6 \pi e^{-2x} + 24 e^{3x}
= 6 (\pi e^{-2x} + 4 e^{3x})
= 6 y.
\]



\item \textit{Utilize uma aproximação linear para estimar o valor de $\sqrt[3]{8,02}$.}

Seja $f(x) = \sqrt[3]{x}$, $x_0 = 8$ e $\Delta x = 0,02$. Então e $f^\prime = \dfrac{1}{3\sqrt[3]{x^2}}$ e tem-se:
\begin{align*}
\sqrt[3]{8,02} = \sqrt[3]{8 + 0,02}
= f(x_0+\Delta x)
& \approx f(x_0) + f^{\prime}(x_0)\Delta x \\
& = \sqrt[3]{8} + \frac{1}{3\sqrt[3]{8^2}}0,02
= 2 + \frac{1}{12}0,02
\approx 2 + 0,0016
= 2,0016.
\end{align*}



\item \textit{Ao esvaziar uma bola de futebol, o seu volume (em $cm^3$) após $t$ segundos é dado por $V(t) = 400 \pi (10 - t)$. Com que velocidade ($cm/s$) o raio diminui quando ele mede $5cm$?}

Como o volume de uma esfera de raio $r$ é $V = V(r) = \frac{4}{3} \pi r^3$, tem-se
\[
\frac{\diff{V}}{\diff{t}}
= \frac{4}{3} \pi 3 r^2 \frac{\diff{r}}{\diff{t}}
= 4 \pi r^2 \frac{\diff{r}}{\diff{t}}
\]

Mas $V(t) = 400 \pi (10 - t)$, então $\frac{\diff{V}}{\diff{t}} = 400 \pi (-1) = -400 \pi$. Disto segue que
\[
\frac{\diff{r}}{\diff{t}}
= \frac{ -400 \pi }{ 4 \pi r^2 }
= \frac{ -100 }{ r^2 }
\]
Em particular, para $r = 5cm$, tem-se $\frac{\diff{r}}{\diff{t}} = -100/25 = -4 cm/s$.

\item \textit{Determine a equação da reta tangente ao gráfico de $f(x) = \dfrac{4}{x + 1}$ em $x = 1$.}

O coeficiente angular $m$ da reta tangente em $(x, f(x))$ é dado pela derivada da função, isto é, $f'(x) = 4 \cdot \dfrac{-1}{(x + 1)^2}$. Então, para $x=1$, tem-se $m=4 \cdot (-1)/2^2 = -1$. Como $f(1) = 4/(1+1) = 2$, a reta tangente tem equação $y - 2 = (-1)(x-1)$, ou seja, $y=-x+3$.

\item \textit{Em que pontos $(x,y)$ da curva $x^\frac{2}{3} + y^\frac{2}{3} = 8$ a reta tangente é paralela à reta $y = x + 4$?}

A reta $y = x + 4$ tem coeficiente angular igual a $1$, então para que ela seja paralela à reta tangente à curva, é preciso que $y^\prime = \frac{\diff{y}}{\diff{x}} = 1$. Por derivação implícita, obtem-se que
\[
\left(x^\frac{2}{3} + y^\frac{2}{3}\right)^\prime = 8^\prime
\Rightarrow \frac{2}{3}x^{\frac{-1}{3}} + \frac{2}{3}y^{\frac{-1}{3}}y^\prime = 0
\]
Então se $y^\prime = 1$, conclui-se que $\frac{2}{3}x^{\frac{-1}{3}} + \frac{2}{3}y^{\frac{-1}{3}} = 0$, ou seja, que $x^{\frac{-1}{3}} = -y^{\frac{-1}{3}}$. Disto conclui-se que $x=-y$.

Substituindo na equação que define a curva, obtem-se
\[
x^\frac{2}{3} + x^\frac{2}{3} = 8
\Rightarrow 2x^\frac{2}{3} = 8
\Rightarrow x^\frac{2}{3} = 4
\Rightarrow x^2 = 4^3 = 64
\Rightarrow x = \pm 8.
\]
Logo, os pontos procurados são $P = (8,-8)$ e $Q=(-8,8)$.

\end{enumerate}

\end{document}
